\section{Magnetic and superconducting parameters}\label{app:E.9}

The magnon band of NiO is set almost entirely by one number, the antiferromagnetic exchange between next-nearest-neighbor nickel ions, and the density-functional values of it collected in one comparison span a factor of two and a half, from the same crystal structure, depending on the functional and on the Hubbard correction used~\cite[Table~III]{fischer2009exchange}. The superconducting gap of aluminum, by contrast, is computed from first principles to within a few percent of the measured \SI{0.18}{meV}~\cite[Table~III]{savrasov1996electron}, but only after one parameter, the Coulomb pseudopotential, is fixed by the measured critical temperature.

\subsection{DFT~$+\,U$ and the Hubbard parameter}

Semilocal functionals spread the $d$ electrons of transition-metal oxides over too many orbitals. In NiO they put nickel $d$ states at the top of the valence band and give a small gap of the wrong character, while DFT~$+\,U$ puts oxygen $p$ states there and gives the charge-transfer gap that experiments measure~\cite[Sect.~V]{cococcioni2005linear}. The correction acts on the occupation matrices $n^{I\sigma}_{mm'} = \sum_{v\vect{k}}f_{v\vect{k}}\langle\psi^\sigma_{v\vect{k}}|P^I_{mm'}|\psi^\sigma_{v\vect{k}}\rangle$ of the localized orbitals of site $I$ with spin $\sigma$, where $P^I_{mm'}$ projects on the atomic orbitals $m$ and $m'$ of the shell~\cite[Eqs.~(5) and~(8)]{cococcioni2005linear}. In its simplest rotationally invariant form it adds
\begin{equation}
  E_U = \frac{U_{\mathrm{H}}}{2}\sum_{I\sigma}\Tr\bigl[n^{I\sigma}(1 - n^{I\sigma})\bigr]
  \label{eq:E.9-U}
\end{equation}
to the energy functional~\cite[Eq.~(9)]{cococcioni2005linear}. In the basis that diagonalizes $n^{I\sigma}$ the trace is $\sum_i\lambda_i(1 - \lambda_i)$ over its eigenvalues $0 \le \lambda_i \le 1$, so the correction penalizes fractional occupations and drives each orbital toward full or empty. Physically, the correction restores the piecewise linearity of the energy as a function of the occupation of the shell, which semilocal functionals violate.

The parameter $U_{\mathrm{H}}$ can be computed rather than fitted. In the linear-response method of Cococcioni and de Gironcoli, a small potential shift $\delta V^U_I$ is applied to the localized orbitals of one site in a supercell, and the response of the occupations $n_J$ is computed twice: once self-consistently, giving the response matrix $R^U_{IJ} = \partial n_I/\partial(\delta V^U_J)$, and once at fixed Kohn--Sham potential, giving the bare response $R^{U,0}_{IJ}$. The Hubbard parameter is the difference of their inverses,
\begin{equation}
  U_{\mathrm{H}} = \bigl[(R^{U,0})^{-1} - (R^U)^{-1}\bigr]_{II} ,
  \label{eq:E.9-Ulr}
\end{equation}
which removes the curvature of the energy that comes from rehybridization of the noninteracting bands~\cite[Eqs.~(15), (18) and~(19)]{cococcioni2005linear}. The convergence parameter is the supercell size, since the perturbation of one site must not interact with its images. Density-functional perturbation theory replaces the supercells by a sum over wavevectors in the primitive cell and reaches the same $U_{\mathrm{H}}$ with better numerical control: for nickel in NiO, the two routes agree to $5\times10^{-3}$\,eV, and the value converges from 7.90 to \SI{8.17}{eV} as the supercell grows from $2\times2\times2$ to $4\times4\times4$ or the equivalent $\vect{q}$ grid is refined~\cite[Table~I]{timrov2018hubbard}.

The value is not unique. It depends on the projectors that define $n^{I\sigma}$, on the functional and on the pseudopotential~\cite[Sect.~V]{timrov2018hubbard}, and on whether $U_{\mathrm{H}}$ is recomputed self-consistently from the DFT~$+\,U$ ground state. For the same nickel shell in NiO, published linear-response values include \SI{4.6}{eV} with nonorthogonal atomic projectors~\cite[Sect.~V]{cococcioni2005linear}, \SI{8.17}{eV} with Löwdin-orthogonalized projectors from a semilocal ground state~\cite[Table~I]{timrov2018hubbard}, and \SI{6.26}{eV} computed self-consistently~\cite[Sect.~II]{dossantos2026comparative}. A $U_{\mathrm{H}}$ is meaningful only together with its projectors, and a calculation that takes $U_{\mathrm{H}}$ from one code and uses it in another changes the physics.

\subsection{Exchange constants}

The constants below are in the convention of the Heisenberg model \cref{eq:7.1-heis}. The literature uses every combination of sign, of single or double counting of pairs, and of spin vectors of length $|\vect{S}|$ or of unit length. The conversions between them are factors of 2 and of $|\vect{S}|^2$~\cite[Table~1]{he2021tb2j}, and a value of $J^{\mathrm{ex}}$ quoted without its Hamiltonian is ambiguous by these factors.

\paragraph{Energy mapping.} The most direct route computes the total energies of several collinear magnetic configurations in DFT~$+\,U$ and equates their differences to those of \cref{eq:7.1-heis} with classical spins. In the fcc lattice of nickel ions in NiO, with $J_1$ between the 12 nearest neighbors and $J_2$ between the 6 next-nearest neighbors, three configurations fix both. In the ferromagnet each ion has all its neighbors parallel. In the antiferromagnet of type I, with ferromagnetic $(001)$ planes alternating in sign, 4 of the 12 nearest neighbors are parallel and 8 antiparallel, and all 6 next-nearest neighbors, at $\pm a$ along the cube axes, lie in planes of the same sign. In the type-II structure of the ground state, with ferromagnetic $(111)$ planes alternating, half of the nearest neighbors are parallel and all next-nearest neighbors antiparallel. Counting each pair once, the energies per ion are
\begin{equation}
  E_{\mathrm{FM}} = (6J_1 + 3J_2)|\vect{S}|^2, \qquad E_{\mathrm{AFI}} = (-2J_1 + 3J_2)|\vect{S}|^2, \qquad E_{\mathrm{AFII}} = -3J_2|\vect{S}|^2 ,
  \label{eq:E.9-map}
\end{equation}
relative to a common reference, and therefore
\begin{equation}
  J_1 = \frac{E_{\mathrm{FM}} - E_{\mathrm{AFI}}}{8|\vect{S}|^2}, \qquad J_2 = \frac{E_{\mathrm{FM}} + 3E_{\mathrm{AFI}} - 4E_{\mathrm{AFII}}}{24|\vect{S}|^2} .
  \label{eq:E.9-J12}
\end{equation}
These are Ref.~\cite[Eqs.~(6) and~(7)]{fischer2009exchange} converted to \cref{eq:7.1-heis}: that reference uses unit vectors, counts each pair twice and puts a minus sign in front, so that its $J_i$ equal $-J_i/2$ of ours for $|\vect{S}| = 1$. The two measured values for NiO that it collects, from Shanker and Singh and from Hutchings and Samuelsen~\cite{hutchings1972measurement}, are $J_2 = 17.3$ and \SI{19.0}{meV} in this convention.

More configurations resolve more neighbors. A variant flips one chosen pair of spins in a large supercell, in the four combinations $\uparrow\uparrow$, $\uparrow\downarrow$, $\downarrow\uparrow$ and $\downarrow\downarrow$, and isolates the coupling of that pair from all others; Ref.~\cite[Sect.~II]{dossantos2026comparative} uses it for NiO in $4\times4\times4$ supercells. In any mapping, the magnetic moment computed in DFT is not $2|\vect{S}|$ Bohr magnetons, and the classical $|\vect{S}|^2$ in \cref{eq:E.9-map} must be replaced by a choice, $|\vect{S}|^2$ or $|\vect{S}|(|\vect{S}| + 1)$, that changes the constants by up to a factor of two for $|\vect{S}| = 1$. The moments may change size from one configuration to the next, which a Heisenberg model cannot describe. The constants also depend on $U_{\mathrm{H}}$: in NiO the superexchange scales roughly as $t_{\mathrm{hop}}^2/U_{\mathrm{H}}$ (\cref{sec:7.1}), so the spread of $U_{\mathrm{H}}$ above becomes a spread of $J_2$, and the density-functional values of Ref.~\cite[Table~III]{fischer2009exchange} range from 11 to \SI{27}{meV} in our convention.

\paragraph{Spin spirals.} A spin spiral of wavevector $\vect{Q}$, with every moment tilted by a cone angle $\theta$ from a fixed axis and rotated in phase by $\vect{Q}\cdot\vect{R}$ from cell to cell, can be computed in the primitive cell when spin--orbit coupling is neglected, by a generalized Bloch theorem. For a lattice with one magnetic ion per cell, \cref{eq:7.1-heis} has $\vect{S}_{\vect{R}}\cdot\vect{S}_0 = |\vect{S}|^2[\sin^2\theta\cos(\vect{Q}\cdot\vect{R}) + \cos^2\theta]$ and therefore
\begin{equation}
  E(\vect{Q},\theta) - E(0,\theta) = \frac{|\vect{S}|^2\sin^2\theta}{2}\bigl[J^{\mathrm{ex}}(\vect{Q}) - J^{\mathrm{ex}}(0)\bigr], \qquad J^{\mathrm{ex}}(\vect{Q}) = \sum_{\vect{R}\neq0}J^{\mathrm{ex}}_{0\vect{R}}\,\ee^{\ii\vect{Q}\cdot\vect{R}} ,
  \label{eq:E.9-spiral}
\end{equation}
per ion, the single-sublattice case of Ref.~\cite[Eq.~(7)]{jacobsson2013exchange}. A set of spirals on a grid of $\vect{Q}$ maps out $J^{\mathrm{ex}}(\vect{Q})$, and a Fourier transform gives the real-space constants out to a range set by the grid. Small cone angles keep the perturbation of the ground state small, which reduces the change of the moments that limits collinear mapping~\cite[Sect.~II]{jacobsson2013exchange}. For a ferromagnet, linear spin-wave theory in the convention of \cref{eq:7.1-heis} has $\omega(\vect{Q}) = |\vect{S}|[J^{\mathrm{ex}}(\vect{Q}) - J^{\mathrm{ex}}(0)]$, so the spiral energy is the magnon energy directly: $\omega(\vect{Q}) = 2[E(\vect{Q},\theta) - E(0,\theta)]/(|\vect{S}|\sin^2\theta)$, the frozen-magnon relation. For a collinear antiferromagnet with two sublattices, the same theory has
\begin{equation}
  \omega(\vect{Q}) = |\vect{S}|\Bigl\{\bigl[J^{\mathrm{ex}}_{\mathrm{AB}}(0) - J^{\mathrm{ex}}_{\mathrm{AA}}(0) + J^{\mathrm{ex}}_{\mathrm{AA}}(\vect{Q})\bigr]^2 - |J^{\mathrm{ex}}_{\mathrm{AB}}(\vect{Q})|^2\Bigr\}^{1/2} ,
  \label{eq:E.9-afm}
\end{equation}
with $J^{\mathrm{ex}}_{\mathrm{AA}}(\vect{Q})$ and $J^{\mathrm{ex}}_{\mathrm{AB}}(\vect{Q})$ the sums of $J^{\mathrm{ex}}_{0\vect{R}}\,\ee^{\ii\vect{Q}\cdot\vect{R}}$ over the sites $\vect{R} \ne 0$ of the same and of the other sublattice; with $J^{\mathrm{ex}}_{\mathrm{AB}} = 0$ it reduces to the ferromagnetic form. The spirals then need one cone angle per sublattice~\cite[Eq.~(8)]{jacobsson2013exchange}, and the two sums follow from their energies.

\paragraph{Green's functions.} The third route needs only the ground state. Rotating two moments by small angles changes the energy, by the magnetic force theorem, by the change of the single-particle energies at fixed potential, which second-order perturbation theory expresses through the Green's function of the ground state. The result is the formula of Liechtenstein, Katsnelson, Antropov and Gubanov, which in a Wannier basis reads
\begin{equation}
  J^{\mathrm{LKAG}}_{ij} = -\frac{1}{4\pi}\int_{-\infty}^{\mu_{\mathrm{ch}}}\dd\varepsilon\,\Imag\Tr\bigl[\Delta_i\,\GF^{\uparrow}_{ij}(\varepsilon)\,\Delta_j\,\GF^{\downarrow}_{ji}(\varepsilon)\bigr]
  \label{eq:E.9-lkag}
\end{equation}
in the convention of its source~\cite[Eq.~(19) and Table~1]{he2021tb2j}, $H = -\sum_{i\neq j}J^{\mathrm{LKAG}}_{ij}\,\hat{\vect{e}}_i\cdot\hat{\vect{e}}_j$ with unit vectors and each pair counted twice, so that $J^{\mathrm{ex}}_{ij} = -2J^{\mathrm{LKAG}}_{ij}/|\vect{S}|^2$ in the convention of \cref{eq:7.1-heis}. Here $\Delta_i$ is the exchange splitting of the on-site Hamiltonian of ion $i$, $\GF^\sigma_{ij}$ the block of the spin-$\sigma$ Green's function between the Wannier orbitals of ions $i$ and $j$, the trace runs over orbitals, and the integral runs up to the chemical potential $\mu_{\mathrm{ch}}$ (\cref{sec:2.5}). One ground-state calculation yields every pair at once, ferromagnetic or antiferromagnetic, and decomposes each constant by orbital. The formula describes infinitesimal rotations about the ground state, while energy mapping measures large ones, and the two differ where the moments are not rigid. For NiO, linear spin-wave theory with constants from energy mapping overestimates the measured magnon energies by about 14 percent, and with constants from the Green's-function formula by up to 40 percent~\cite[Sect.~III]{dossantos2026comparative}. The convergence parameters of \cref{eq:E.9-lkag} are the Wannier basis, which must contain the magnetic $d$ orbitals and the ligand $p$ orbitals that mediate superexchange, and the $\vect{k}$ grid of the Green's function, which sets the range of $|\vect{R}_i - \vect{R}_j|$ that is free of images.

\paragraph{Spin response.} The most complete route does without the Heisenberg model. Time-dependent density-functional perturbation theory computes the transverse spin susceptibility of the DFT~$+\,U$ ground state as a function of $(\vect{Q},\omega)$, and its poles are the magnons; constants are then fitted to the dispersion only for comparison~\cite[Sect.~II]{dossantos2026comparative}. This is the spin analogue of the dielectric response of \cref{app:E.6}, and it costs correspondingly more.

\subsection{Eliashberg theory}

The gap and critical temperature of a phonon-mediated superconductor follow from the Eliashberg equations, the Nambu--Gor'kov form of the electron self-energy with the phonon propagator of \cref{sec:6.4} and the Coulomb repulsion. On the imaginary axis, at the fermionic Matsubara frequencies $\omega^{\mathrm{F}}_j$ of \cref{sec:4.7}, and for an isotropic gap restricted to the Fermi surface, they are two coupled equations for a mass-renormalization function $Z^{\mathrm{El}}(\ii\omega^{\mathrm{F}}_j)$ and a gap function $\Delta^{\mathrm{El}}(\ii\omega^{\mathrm{F}}_j)$,
\begin{align}
  Z^{\mathrm{El}}(\ii\omega^{\mathrm{F}}_j) &= 1 + \frac{\pi T}{\omega^{\mathrm{F}}_j}\sum_{j'}\frac{\omega^{\mathrm{F}}_{j'}}{\sqrt{(\omega^{\mathrm{F}}_{j'})^2 + \Delta^{\mathrm{El}}(\ii\omega^{\mathrm{F}}_{j'})^2}}\;\lambda_{\mathrm{ep}}(\omega^{\mathrm{F}}_j - \omega^{\mathrm{F}}_{j'}) , \label{eq:E.9-Z}\\
  Z^{\mathrm{El}}(\ii\omega^{\mathrm{F}}_j)\,\Delta^{\mathrm{El}}(\ii\omega^{\mathrm{F}}_j) &= \pi T\sum_{j'}\frac{\Delta^{\mathrm{El}}(\ii\omega^{\mathrm{F}}_{j'})}{\sqrt{(\omega^{\mathrm{F}}_{j'})^2 + \Delta^{\mathrm{El}}(\ii\omega^{\mathrm{F}}_{j'})^2}}\Bigl[\lambda_{\mathrm{ep}}(\omega^{\mathrm{F}}_j - \omega^{\mathrm{F}}_{j'}) - \mu^*_{\mathrm{C}}\Bigr] , \label{eq:E.9-Delta}
\end{align}
with the frequency-dependent coupling
\begin{equation}
  \lambda_{\mathrm{ep}}(\omega^{\mathrm{B}}) = \int_0^\infty\dd\omega\,\frac{2\omega\,\alpha^2F(\omega)}{\omega^2 + (\omega^{\mathrm{B}})^2} ,
  \label{eq:E.9-lambda}
\end{equation}
evaluated at the bosonic differences $\omega^{\mathrm{B}} = \omega^{\mathrm{F}}_j - \omega^{\mathrm{F}}_{j'}$, so that $\lambda_{\mathrm{ep}}(0)$ is the coupling constant $\lambda_{\mathrm{ep}}$ of \cref{sec:6.4}. These are the isotropic forms of the Fermi-surface-restricted equations of Ref.~\cite[Eqs.~(41), (42) and~(48)]{lee2023electron}; the literature writes $Z$, $\Delta$, $\lambda$ and $\mu^*$, symbols that this book uses for other things. The highest temperature at which \cref{eq:E.9-Delta} has a nonzero solution is $T_c$, and the gap on the real axis follows by analytic continuation of $\Delta^{\mathrm{El}}(\ii\omega^{\mathrm{F}}_j)$, usually with Padé approximants (\cref{app:C.5}).

The Coulomb pseudopotential $\mu^*_{\mathrm{C}}$ is the weak point. It replaces the screened Coulomb repulsion between paired electrons, renormalized down from its bare Fermi-surface value by the retardation of the phonon interaction, and in most calculations it is chosen, with values from 0.1 to 0.2 giving reasonable agreement with experiment~\cite[Sect.~4]{lee2023electron}. Since $T_c$ depends on it exponentially, a computed $T_c$ with a chosen $\mu^*_{\mathrm{C}}$ is a consistency check on $\alpha^2F$ rather than a prediction. Two alternatives remove the parameter: computing $\mu^*_{\mathrm{C}}$ as a Fermi-surface average of the screened interaction, and density-functional theory for superconductors~\cite[Sect.~4]{lee2023electron}.

A quick estimate replaces the equations by the McMillan formula in the form of Allen and Dynes,
\begin{equation}
  T_c = \frac{\omega_{\log}}{1.2}\exp\left[-\frac{1.04(1 + \lambda_{\mathrm{ep}})}{\lambda_{\mathrm{ep}} - \mu^*_{\mathrm{C}}(1 + 0.62\lambda_{\mathrm{ep}})}\right], \qquad
  \omega_{\log} = \exp\left[\frac{2}{\lambda_{\mathrm{ep}}}\int_0^\infty\dd\omega\,\frac{\alpha^2F(\omega)}{\omega}\ln\omega\right] ,
  \label{eq:E.9-AD}
\end{equation}
which is Ref.~\cite[Eq.~(30)]{lee2023electron}, a fit to numerical solutions of the Eliashberg equations useful for preliminary estimates.

\paragraph{Convergence and cost.} The input $\alpha^2F$ is a double Fermi-surface average of $|g|^2$, with two energy delta functions, and it converges only on fine grids of $\vect{k}$ and $\vect{Q}$, which in practice means Wannier interpolation of the vertices (\cref{app:E.5}). The anisotropic equations, with a gap that varies over the Fermi surface, need the vertices on commensurate fine grids, restricted to an energy window of a few times the largest phonon energy around the Fermi level~\cite[Sect.~4]{lee2023electron}. The Matsubara sum needs frequencies up to a cutoff, commonly ten times the largest phonon frequency, with convergence tests recommended~\cite[Sect.~4]{lee2023electron}. Since the spacing is $2\pi T$, the number of frequencies grows as $1/T$, and low-$T_c$ superconductors are expensive: for aluminum, with phonons up to \SI{40.5}{meV} ($\omega_{\mathrm{max}} \approx \SI{470}{K}$~\cite[Sect.~III]{savrasov1996electron}), a cutoff of ten times that holds about 750 positive frequencies at \SI{1}{K} and 7500 at \SI{0.1}{K}. Sparse sampling of the Matsubara axis reduces the count~\cite[Sect.~4]{lee2023electron}. The analytic continuation to the real axis is the last numerical hazard: Padé approximants can produce spurious structure near the gap edge, where the pair-breaking absorption of \cref{chap:15} begins, and a real-axis solution or a check against the imaginary-axis data is needed there.

\begin{workedexample}{Aluminum from its Eliashberg function}\label{we:E-al}
Linear-response calculations of aluminum give $\lambda_{\mathrm{ep}} = 0.44$, against 0.42 from tunneling, and a logarithmic phonon frequency $\omega_{\log} = \SI{270}{K}$~\cite[Tables~II and~III]{savrasov1996electron}. Solving the Eliashberg equations with this $\alpha^2F$, the measured $T_c = \SI{1.18}{K}$ is reproduced with $\mu^*_{\mathrm{C}} = 0.12$, and the zero-temperature gap is then \SI{0.18}{meV}, equal to the measured value~\cite[Table~III]{savrasov1996electron}. The weak-coupling ratio of \cref{sec:7.2} agrees, $\Delta_{\mathrm{BCS}} = 1.764\,T_c = \SI{0.179}{meV}$, as it should for $\lambda_{\mathrm{ep}} = 0.44$.

\Cref{eq:E.9-AD} with the same three numbers gives $T_c = (270/1.2)\,\ee^{-5.213}\,\si{K} = \SI{1.22}{K}$, within four percent of the full solution. Its sensitivity to $\mu^*_{\mathrm{C}}$ is the logarithmic derivative of the exponent, $\partial\ln T_c/\partial\mu^*_{\mathrm{C}} = -1.04(1 + \lambda_{\mathrm{ep}})(1 + 0.62\lambda_{\mathrm{ep}})/[\lambda_{\mathrm{ep}} - \mu^*_{\mathrm{C}}(1 + 0.62\lambda_{\mathrm{ep}})]^2 = -23$. A change of $\mu^*_{\mathrm{C}}$ by $\pm0.02$, well inside the range used in the literature, moves $T_c$ by a factor of 0.60 or 1.53, from \SI{0.74}{K} at $\mu^*_{\mathrm{C}} = 0.14$ to \SI{1.87}{K} at 0.10. The gap and $T_c$ of aluminum are measured, and the first-principles $\alpha^2F$ enters where a measured quantity is not available, in the frequency dependence of the phonon-assisted pair-breaking absorption of \cref{chap:15}.
\end{workedexample}

Wherever a magnetic or superconducting target is ranked by a computed parameter, the spread over $U_{\mathrm{H}}$ and mapping, or over $\mu^*_{\mathrm{C}}$, is the theory uncertainty (\cref{sec:8.8}).
